% ============================================================================
% Chapter 5: Standards and Design Constraints  (FYDP-II)
%
% Replaces the whole FYDP-I Chapter 5 (its 5.1 complex-engineering text and
% Table 5.2 are rewritten here as 5.4 and 5.3). Compiles with pdfLaTeX.
% Uses \checkmark (amssymb); the line below gives a fallback if it is missing.
% Citation keys: banglabert, mmbert, gemma4, llamacpp (new_references.tex).
% ============================================================================

\providecommand{\checkmark}{$\surd$}

\chapter{Standards and Design Constraints}
\label{ch:standards}

This chapter describes the standards the system follows, the constraints that
shaped its design, its cost, and how the project meets the criteria of a
complex engineering problem. Figures come from the measurements in
Chapter~\ref{ch:implementation} unless marked as estimates.

% ----------------------------------------------------------------------------
\section{Compliance with the Standards}
\label{sec:standards}

The system is software that runs on commodity hardware and exchanges data over
standard web protocols. It was built on published, widely adopted standards
rather than vendor-specific technology, so that it can be installed, operated
and integrated without special equipment or proprietary services.

\subsection{Software Standards}

\begin{itemize}
  \item \textbf{Unicode and UTF-8 (RFC 3629).} All Bangla text is stored,
  processed and exchanged as Unicode in UTF-8 encoding. The word segmentation
  rule is defined on the Unicode Bengali block (U+0980 to U+09FF) and on the
  Bangla full stop (U+0964), which the Unicode standard shares with other
  Indic scripts. This keeps the text readable by any standard tool.
  \item \textbf{Data formats.} The dataset is a CSV file, word labels are stored
  as JSON (RFC 8259), and results, logs and LLM-judge responses are written as
  JSON and JSON Lines. Trained models are saved in the open
  \texttt{safetensors} format, which, unlike Python pickle files, cannot execute
  code when loaded.
  \item \textbf{Requirements specification.} The functional and non-functional
  requirements in Section~3.1.1 follow the structure of ISO/IEC/IEEE 29148 for
  requirements engineering, and each requirement is checked against a
  measurement in Table~\ref{tab:requirements}.
  \item \textbf{Software quality.} The design addresses the quality
  characteristics of the ISO/IEC 25010 product quality model: performance
  efficiency (57~ms per check on a CPU), reliability (input validation and the
  tests in Table~\ref{tab:tests}), maintainability (modular code with one
  configuration file), portability (runs on CPU or any DirectX~12 GPU) and
  security (no data leaves the machine). The project was not formally assessed
  against the standard.
  \item \textbf{Coding conventions and version control.} The Python code follows
  PEP~8 conventions, dependencies are listed with version constraints, and all
  code, results and documentation are kept under Git version control.
  \item \textbf{Open-source licences.} Every component is used within its
  licence: Apache~2.0 (Gemma~4 models, Transformers, PEFT, Streamlit), MIT
  (mmBERT, llama.cpp, FastAPI), BSD (PyTorch, scikit-learn, pandas, Uvicorn),
  and CC BY-NC-SA 4.0 for BanglaBERT, which restricts it to non-commercial use.
\end{itemize}

\subsection{Hardware Standards}

\begin{itemize}
  \item \textbf{x86-64 architecture.} The system runs on a standard x86-64
  desktop (AMD Ryzen 5 7500F, AM5 platform, DDR5 memory) under Windows 11.
  \item \textbf{PCI Express 4.0.} The AMD Radeon RX 6600 GPU connects over the
  standard PCI Express 4.0 interface; no special hardware is needed.
  \item \textbf{DirectX 12 and DirectML.} GPU training and inference use
  Microsoft DirectML, which runs on any DirectX~12 capable GPU from AMD, Intel
  or NVIDIA. This avoids depending on one vendor's platform, such as NVIDIA
  CUDA, which the planned hardware in FYDP-I would have required.
  \item \textbf{Vulkan.} The LLM judges were run through the Vulkan graphics and
  compute API, an open standard of the Khronos Group supported by all major GPU
  vendors.
  \item \textbf{CPU-only operation.} The deployed model does not need a GPU at
  all: it meets the latency requirement on the CPU alone, so it runs on any
  ordinary computer with 16~GB of memory.
\end{itemize}

\subsection{Communication Standards}

\begin{itemize}
  \item \textbf{HTTP.} The backend is served by Uvicorn, an implementation of
  the ASGI server interface, over HTTP/1.1 (RFC 9110 and RFC 9112).
  \item \textbf{REST and JSON.} The FastAPI backend exposes a REST interface: a
  client sends a \texttt{POST /detect} request whose JSON body holds the
  context, question and answer, and receives a JSON response with the verdict,
  the highlighted spans and a probability for every word. \texttt{GET /health}
  reports whether a model is loaded. Invalid requests are rejected with the
  standard HTTP status code 422.
  \item \textbf{OpenAPI.} FastAPI publishes an OpenAPI~3 description of the
  interface, so other systems, such as a RAG pipeline, can generate client
  code for it and read its documentation in a browser.
  \item \textbf{WebSocket.} The Streamlit dashboard communicates with the
  browser over the WebSocket protocol (RFC 6455).
  \item \textbf{Local communication.} By default both services listen only on
  the local machine (127.0.0.1). The LLM judges were also called through a
  local HTTP interface compatible with the widely used chat-completions
  format, so no text was sent over the internet at any point.
\end{itemize}

% ----------------------------------------------------------------------------
\section{Design Constraints}
\label{sec:constraints}

This section describes the constraints that shaped the design and how the
implemented system meets them.

\subsection{Economic Constraint}

The project had no budget for cloud computing or paid APIs, so every component
had to run on hardware the team already owned and use free software. This was
met in full. All training, evaluation and inference ran on one consumer desktop
with an AMD RX 6600 GPU; the Google Colab subscription planned in FYDP-I was not
needed. All libraries are open source, and the pretrained models are free to
download.

The running cost of the finished system is close to zero. Checking an answer
takes 57~ms on an ordinary CPU and needs no GPU, no internet connection and no
per-request fee. The alternative the project set out to replace, an LLM used as
a judge, is far more expensive to run: on the same GPU, a local
12-billion-parameter judge needed 25.7 minutes to check the 588 test answers, where
BanglaBERT needed about 7 seconds, and a cloud LLM API would add a charge for
every request.

One economic limit concerns licensing. BanglaBERT is released under the Creative
Commons BY-NC-SA 4.0 licence~\cite{banglabert}, which allows research and
educational use but not commercial use. A commercial product would need
permission from its authors, or a backbone with a permissive licence. The mmBERT
comparison model is MIT-licensed~\cite{mmbert} and could be used commercially,
but it was slightly less accurate and too slow on a CPU (210~ms per answer).

\subsection{Environmental Constraint}

The system had to keep energy use low, both to train and, more importantly, in
everyday use, where it may check thousands of answers. A small model meets this
directly: one check uses 11.5~ms of GPU time, against 2,560~ms for the 12B LLM
judge, so the same work costs roughly 220 times less computation. On a CPU the
model needs no GPU at all.

Training was also modest. All experiments together, including the three-seed
runs, the comparison models and the LLM judges, took about 8 hours of machine
time on one desktop. At typical power draw for such a machine this is roughly
1 to 2~kWh (an estimate, not a measurement). LoRA reduced GPU training time per
epoch by a factor of 1.8. No new hardware was bought and no data-centre
resources were used, which avoids the embodied carbon of dedicated equipment.

\subsection{Ethical Constraint}

\paragraph{Privacy.} RAG systems often answer questions about health,
government services or personal matters. Because the detector runs entirely on
the local machine, the question, the retrieved context and the answer never
leave the device. No third-party service receives user data, unlike an
approach based on a cloud LLM API.

\paragraph{Transparency.} The system does not give an unexplained verdict. It
highlights the exact words it considers unsupported and shows a probability
for each word, so a user can check its reasoning against the context. The demo
also shows the lexical baseline beside the model.

\paragraph{Honest limits.} The detector is not perfect: it misses about 11\% of
hallucinated answers and wrongly flags about 4\% of faithful ones. It must be
presented as an aid to human judgement, not as a guarantee of correctness. The
results in this report are given with their baselines, with the variation
across three training runs, and with the system's weak types stated openly.

\paragraph{Data.} The dataset was built from publicly available sources: Bangla
Wikipedia (licensed CC BY-SA 4.0) and public Bangladeshi government portals. It
contains no personal data about private individuals. The pretrained models are
used within their licence terms, which for BanglaBERT means non-commercial
research use. Human annotation of the dataset is done by the team members
themselves.

\subsection{Health and Safety Constraint}

The system has no physical hazards; it is software running on standard office
computers. Its safety concern is informational. 500 rows of the dataset come
from the healthcare domain, including material from the Directorate General of
Health Services. A hallucinated medical answer that the detector fails to flag
could lead a user to act on wrong information, while a false alarm could make
a user doubt correct advice.

The design responds to this in three ways. The detector reports which words
are unsupported, so a reader can verify them rather than trust a single score.
The results are broken down by domain and hallucination type, so weak areas are
known before deployment. And the system is intended to support, not replace,
human review: in a health setting its output should be shown with a clear
notice that flagged and unflagged answers alike must be checked against an
authoritative source.

\subsection{Social Constraint}

Bangla is one of the most widely spoken languages in the world, yet it has few
natural language processing tools, and existing hallucination detectors and
datasets are almost all English. This project brings hallucination detection
to Bangla RAG systems, so that Bangla-speaking users of chatbots and question
answering services can be warned about unsupported answers.

The system was also designed to be affordable to the institutions most likely
to deploy it in Bangladesh, such as universities, small companies and
government offices, which often lack GPUs and budgets for cloud AI services.
It runs on an ordinary computer without a GPU. Highlighting the wrong words,
rather than giving a bare verdict, makes the output understandable to users
without technical training. A social risk remains: users may trust any answer
the detector does not flag. This is why the system presents itself as a
warning tool with known error rates.

\subsection{Political Constraint}

Two of the eight domains, Bangladesh Affairs and Government Services, include
political and administrative content such as elections, public office holders
and government procedures. Hallucinations on these subjects can spread
misinformation about public institutions. The detector is politically neutral
in one precise sense: it judges an answer only against the retrieved context
and has no opinion of its own on any political question. The consequence is
that it inherits the reliability of its source. If the retrieved context is
biased or wrong, an answer that repeats it will be judged supported. Deployers
must therefore choose trustworthy sources for retrieval.

Running locally also supports data sovereignty. Government bodies can check
answers without sending citizens' questions or internal documents to foreign
cloud providers, and the system keeps working without dependence on any
external service.

\subsection{Sustainability}

The system is designed to stay useful and maintainable after the project ends.
\begin{itemize}
  \item \textbf{Cheap to retrain.} A full fine-tuning run takes 25 to 55 minutes
  on the RX 6600, and a LoRA run 13 to 29 minutes, so the model can be retrained
  whenever the dataset grows or is corrected by human annotation.
  \item \textbf{Cheap to extend.} Each LoRA adapter is 3.4~MB, so separate
  adapters for new domains can be added without storing another full model.
  \item \textbf{Open components.} All software is open source and all models
  have public weights, so the system does not depend on a vendor that could
  change its prices or withdraw its service.
  \item \textbf{Robust to data changes.} The dataset is still being extended.
  The pipeline measures sequence lengths and checks the data format on every
  load rather than assuming fixed dimensions, so new versions can be used
  without code changes.
\end{itemize}

Two risks to long-term maintenance remain. GPU training relies on DirectML,
which Microsoft now maintains only lightly and which requires an older PyTorch
release; a future move to a supported GPU platform may be needed. And
BanglaBERT reads at most 512 tokens, which may become a limit if future RAG
contexts are much longer.

% ----------------------------------------------------------------------------
\section{Cost Analysis}
\label{sec:cost}

Table~\ref{tab:cost} gives the actual cost of FYDP-II. It replaces the estimate
in the FYDP-I report.

\begin{table}[htbp]
\centering
\small
\caption{Project cost analysis (actual, FYDP-II)}
\label{tab:cost}
\begin{tabular}{p{0.4cm}p{3.2cm}p{5.6cm}p{2.2cm}}
\hline
\textbf{No.} & \textbf{Item} & \textbf{Description} & \textbf{Cost (USD)} \\
\hline
1 & Development hardware & Existing laptops and one desktop with an AMD RX 6600 GPU & 0 (existing) \\
2 & Cloud GPU & Planned Google Colab Pro (about \$20) was not needed; all training ran locally & 0 \\
3 & Software & Python, PyTorch, Transformers, PEFT, FastAPI, Streamlit, llama.cpp & 0 (open source) \\
4 & Pretrained models & BanglaBERT (CC BY-NC-SA 4.0), mmBERT (MIT), Gemma~4 judges (Apache 2.0) & 0 \\
5 & Dataset & Built by the team from public sources & 0 \\
6 & LLM-judge comparison & Run locally instead of through a paid API & 0 \\
7 & Electricity & About 8 hours of machine time, roughly 1--2 kWh (estimate) & under 1 \\
8 & Documentation & Report printing, binding and presentation material & about 15 \\
\hline
 & \textbf{Total} & & \textbf{about 16} \\
\hline
\end{tabular}
\end{table}

The actual cost is below the FYDP-I estimate of about \$35, mainly because GPU
training on local hardware made the cloud subscription unnecessary. The cost of
using the system is also far below that of an LLM judge: it needs no API
subscription, and it checks an answer about 220 times faster than a local 12B
LLM on the same GPU.

\paragraph{Cost of operation.} A deployment needs one of two files: the fully
fine-tuned model (420~MB), or a LoRA adapter (3.4~MB) loaded on top of the
public BanglaBERT base model. Any computer with 16~GB of memory can serve it
through the FastAPI backend, with no licence fee, subscription or per-request
charge. At batch size 1 it checks about 17 answers per second on the CPU used
here and about 87 per second on the RX 6600 GPU. The only running cost is the
electricity of a computer that is already switched on.

% ----------------------------------------------------------------------------
\section{Complex Engineering Problem}
\label{sec:complex}

The Washington Accord defines complex engineering problems by seven attributes
(P1 to P7) and complex engineering activities by five (A1 to A5). This section
shows how the project meets them.

\subsection{Complex Problem Solving}

\begin{itemize}
  \item \textbf{P1: Depth of knowledge required.} The project needs knowledge of
  transformer encoders, token classification, subword tokenisation of an
  agglutinative language, parameter-efficient fine-tuning with LoRA, and
  evaluation methods for sequence labelling. It also needed knowledge of GPU
  computing to train on DirectML.
  \item \textbf{P2: Range of conflicting requirements.} High accuracy conflicts
  with low latency, small memory use and running on a CPU. Larger models such as
  mmBERT or an LLM judge were either slower or less accurate; the design had to
  find a model that meets all requirements at once.
  \item \textbf{P3: Depth of analysis required.} The project had to establish
  whether the model truly learns the task. That required a baseline that shows
  how far word overlap alone gets, a split by context that prevents leakage,
  three training seeds to measure variation, word- and span-level metrics, and
  an analysis by hallucination type.
  \item \textbf{P4: Familiarity of issues.} Hallucination detection for Bangla
  has no established method, dataset or benchmark. Issues such as Bangla case
  suffixes defeating word matching, or the danda attached to words, had to be
  discovered and solved during the project.
  \item \textbf{P5: Extent of applicable codes.} No standard or code of practice
  covers evaluating hallucination detectors, let alone for Bangla. The project
  had to define its own evaluation protocol: the metrics, the split, the
  baselines and the latency measurement method.
  \item \textbf{P6: Extent of stakeholder involvement.} The system serves
  developers who integrate RAG into applications, organisations that must keep
  data private, and end users who rely on answers in areas such as health and
  government services. Their needs for accuracy, privacy and low cost conflict
  and had to be balanced.
  \item \textbf{P7: Interdependence.} The system only works if every part fits
  the others: data validation, word segmentation, subword alignment, the
  encoder and its classifier, span aggregation, the API and the user interface.
  A mismatch at any step, such as the answer being split into different words
  at inference than in training, silently breaks the result.
\end{itemize}

\begin{table}[htbp]
\centering
\caption{Mapping with complex problem solving}
\label{tab:complex-problems}
\begin{tabular}{llc}
\hline
\textbf{Attribute} & \textbf{Description} & \textbf{Addressed} \\
\hline
P1 & Depth of knowledge required & \checkmark \\
P2 & Range of conflicting requirements & \checkmark \\
P3 & Depth of analysis required & \checkmark \\
P4 & Familiarity of issues & \checkmark \\
P5 & Extent of applicable codes & \checkmark \\
P6 & Extent of stakeholder involvement & \checkmark \\
P7 & Interdependence & \checkmark \\
\hline
\end{tabular}
\end{table}

\subsection{Engineering Activities}

\begin{itemize}
  \item \textbf{A1: Range of resources.} The project combined people (a
  four-member team sharing data, modelling, evaluation and interface work),
  hardware (a CPU and an AMD GPU), data (a 4,000-row dataset built from
  Wikipedia and government portals), pretrained models (BanglaBERT, mmBERT,
  Gemma~4) and a wide software stack.
  \item \textbf{A2: Level of interaction.} Technical decisions interacted
  strongly. The choice of model set the sequence length, which set the
  latency and memory use, which set the hardware and batch size; the choice of
  licence limited commercial use; the GPU platform forced changes to the model
  code. Each had to be resolved against the others.
  \item \textbf{A3: Innovation.} To our knowledge this is the first token-level
  hallucination detector for Bangla RAG answers. The project also measured the
  cost of LLM judges against a small model on the same data and hardware, and
  adapted model training to run on an AMD consumer GPU through DirectML.
  \item \textbf{A4: Consequences for society and the environment.} Undetected
  hallucinations in health or government answers can mislead people, and false
  alarms can erode trust. Running locally protects privacy, and a small model
  uses about 220 times less computation per check than an LLM judge.
  \item \textbf{A5: Familiarity.} The work went beyond the team's previous
  experience: a new language setting, a new hardware platform and a new
  evaluation protocol. Each was approached from first principles and verified
  by measurement.
\end{itemize}

\begin{table}[htbp]
\centering
\caption{Mapping with complex engineering activities}
\label{tab:complex-activities}
\begin{tabular}{llc}
\hline
\textbf{Attribute} & \textbf{Description} & \textbf{Addressed} \\
\hline
A1 & Range of resources & \checkmark \\
A2 & Level of interaction & \checkmark \\
A3 & Innovation & \checkmark \\
A4 & Consequences for society and the environment & \checkmark \\
A5 & Familiarity & \checkmark \\
\hline
\end{tabular}
\end{table}

% ----------------------------------------------------------------------------
\section{Summary}
\label{sec:ch5summary}

The system is built on open standards: Unicode text, standard data and model
formats, standard GPU and web interfaces, and a REST API described by OpenAPI.
It meets its design constraints: it costs nothing to run, uses little energy,
keeps data on the local machine, and serves Bangla users on ordinary hardware.
Its main limits are the non-commercial licence of BanglaBERT and the need for
human oversight in sensitive domains. The whole project cost about \$16. It
meets all seven attributes of a complex engineering problem and all five of
complex engineering activities.
